\documentclass[11pt]{article}
\usepackage[a4paper,margin=26mm]{geometry}
\usepackage{amsmath,amssymb,amsthm,mathrsfs,hyperref}
\setlength{\emergencystretch}{3em}
\newtheorem{theorem}{Theorem}
\newtheorem{lemma}[theorem]{Lemma}
\newtheorem{proposition}[theorem]{Proposition}
\newcommand{\MA}{\operatorname{MA}}
\newcommand{\PSH}{\operatorname{PSH}}
\newcommand{\vol}{\operatorname{vol}}
\newcommand{\EE}{\mathcal E^1_\theta}
\title{Energy-complete unbounded moving systems\protect\\
and their paired relative weighted limits}
\author{Unreviewed conditional research, continuation 203}
\date{6 October 2026}
\begin{document}
\maketitle
\begin{abstract}
In a big rational class, truncation reduces simultaneous unbounded
full-system limits to uniformly bounded generated quotients.
A finite absolute moment estimate controls the energy difference of
the truncated systems, without any quadratic estimate for unbounded
models. Finite strict locality controls the total variation lost by
truncation. Under explicitly unresolved genuine scalar and finite-list
moving foundations, this identifies every energy-complete joint
diagonal and mixed measure, independently of degree selection.
The same finite absolute estimate identifies paired moving relative
tests, even when their canonical indices change independently.
This does not prove weighted convergence of a fixed original
unbounded function against every physical auxiliary measure:
the singular-reference discrepancy for that assertion is retained.
All results remain unreviewed and conditional. Original P16.5 is open.
\end{abstract}

\section{Whole domain and exact unresolved inputs}
Let $X$ be a smooth integral projective $n$-fold, $n\geq1$, over
any algebraically closed trivially valued characteristic-zero field.
Let $L$ be any actual big rational line bundle, $\theta=c_1(L)$,
$K=X^{\mathrm{an}}$, and $V=\vol(L)>0$.
Keep every supported semivaluation and every central multiplicity.
The genuine cone is the all-point decreasing limit of one FS net on
varying actual rational bundles whose numerical classes tend to
$\theta$, with finite generic endpoint. Normalize at the generic
trivial valuation $v_0$; put $R=P_\theta(0)$.
Choose an actual Kodaira identity $L=aA+E_0$, with $A$ ample,
$a>0$ rational and $E_0$ effective, and put
$Z=E_0^{\mathrm{an}}$, $U=K\setminus Z$.
No field countability, finite generation, invariant metric or support
restriction is added.

The paper assumes the following explicit \emph{conditional} inputs.
They are not asserted to hold on the actual big domain.
\begin{enumerate}
\item[(S)] The literal genuine cone is convex and decreasing/max closed,
has all-point envelopes and reference $R$, and has the whole concave,
monotone, decreasing-net continuous, $V$-translation equivariant
outer full-section energy $E$, with $E(R)=0$.
Its finite-energy space is the actual set
$\EE=\{u\in\PSH(\theta):E(u)>-\infty\}$.
Every normalized $u\in\EE$ has nonpositive rational model obstacles
$f_i\downarrow u$, with $f_i(v_0)=0$ and
$u_i=P_\theta(f_i)\downarrow u$ all-point, and
$E(u_i)\downarrow E(u)$.
\item[(M)] The original target-independent positive Radon map
$M:\EE\to\mathcal M^+(K)$ has raw mass $V$, constant invariance,
finite-model compatibility, genuine decreasing weak continuity
and $M(u)(Z)=0$.
The positive genuine mixed map has diagonal $M$, and satisfies
the finite polarization identities and
\[
 0\leq M(u_1,\ldots,u_n)
 \leq \frac{n^n}{n!}M\!\left(\frac1n\sum_{a=1}^n u_a\right).
                                                        \tag{1}
\]
Its canonical weighted continuity and all genuine relative-test
integrability are retained.
\item[(F)] Resolve the full $|mL|$ system and any fixed finite list of
full obstacle lattices on common smooth equivariant models:
\[
 p_m^*L=A_m+T_m,\qquad mA_m\text{ generated},\quad mT_m\geq0,
 \qquad V_m=A_m^n\longrightarrow V.                      \tag{2}
\]
For a nonpositive model obstacle $f$, the relatively generated
auxiliary scalar $\psi_{m,f}\leq0$ has physical stage
$q_{T_m}+\psi_{m,f}\leq P_\theta(f)$ on every lift.
Its energy $e_m(\psi_{m,f})$ tends to $E(P_\theta(f))$,
its generic value tends to zero, and its pushed top measure tends
to $M(P_\theta(f))$.
At a fixed finite list these statements hold on one common actual
$A_m$, and full filtration order gives
$\psi_{m,f}\leq\psi_{m,g}$ whenever $f\leq g$.
 For fixed $f$, the quotient is uniformly bounded and converges
 uniformly on each compact subset of $U$ to $P_\theta(f)-R$.
 Both $R$ and these model envelopes are continuous on $U$, and
 $q_{T_m}\to R$ uniformly on every compact subset of $U$.
 All degree limits are divisibility-cofinal with nonempty systems.
\item[(B)] A uniformly bounded relatively generated quotient on
the data (2), with generic error tending to zero, physical stage
below a fixed $w\in\EE$ continuous on $U$, and compact-uniform
quotient convergence to $w-R$, has auxiliary energy tending to
$E(w)$ and pushed raw measure tending to $M(w)$.
This applies also to rational convex combinations of a fixed finite
list of model envelopes and to their fixed relative truncations.
\end{enumerate}
These are the explicitly conditional scalar/moving package of
the whole round159 paper and the conditional conclusions of
continuations201--202, not caller theorem grants. Their complete
proof bodies have been read. The only caller NA grants remain
G0 in its actual ample domain and G1 for actual relative generation
and one FS net. Proving all actual antecedents is separate work.

\section{Finite estimates, with generic terms retained}
On a generated actual rational nef $A$ write $V_A=A^n$,
$e_A(0)=0$, and $\mu_\xi=\MA_{A,\mathrm{raw}}(\xi)$ for a
bounded relatively generated model scalar $\xi$.
All these measures include their central multiplicities.
For $\xi\leq0$ let $B_A(\xi)=-e_A(\xi)$ and
$J_A(\xi)=V_A\xi(v_0)-e_A(\xi)$.
Thus $0\leq J_A(\xi)\leq B_A(\xi)$.
All finite formulas below are at fixed models; they assume no
unbounded auxiliary NA operator.

We use the reconstructed finite absolute estimate of round159:
with $\alpha=2^{-n}$ and
$\delta_A(\xi,\eta)=e_A((\xi+\eta)/2)-(e_A(\xi)+e_A(\eta))/2$,
\[
 \int|\xi-\eta|\,d\mu_\zeta
 \leq V_A|\xi(v_0)-\eta(v_0)|
   +C_n\delta_A(\xi,\eta)^\alpha
       J_A(\zeta)^{1/2}
       \max(J_A(\xi),J_A(\eta),J_A(\zeta))^{1/2-\alpha}.
                                                        \tag{3}
\]
For $n=1$ the exponent-zero factor is omitted.
Its graph proof uses $-D^2\Lambda\geq0$ with every actual
multiplicity, contraction under $D\mapsto|D|$, mixed
Cauchy--Schwarz, and the correctly oriented $n-j$ energy tangent.
This finite assertion is unreviewed proof data here.
If $\xi\geq\eta$ and all three nonpositive costs are at most $D$,
then monotonicity bounds
$0\leq\delta_A(\xi,\eta)\leq(e_A(\xi)-e_A(\eta))/2$,
so (3) gives
\[
 \int(\xi-\eta)\,d\mu_\zeta
 \leq V_A|g_\xi-g_\eta|
      +K_nD^{1-\alpha}(e_A(\xi)-e_A(\eta))^\alpha ,
 \qquad g_\xi=\xi(v_0).                                  \tag{4}
\]
The constants do not depend on model amplitudes, refinements or
positive raw volume. No division by $V_A$ is made.

Finite bounded nef locality is obtained from G0's finite ample
locality by adding one common actual $\varepsilon H$ ample
zero metric, retaining every component coefficient, and taking
the finite polynomial limit $\varepsilon\downarrow0$.
For models between $-C$ and zero, a fixed time twist $c>C+1$
makes their resolved compactified bundles globally generated:
retain only generator weights in $[-mC,0]$ and the weight-zero
generating system of $mA_m$.
The actual nef bundles satisfy
$\Gamma\leq\gamma=p^*L+cF$, with $F$ the full fiber.
Writing any fixed original vertical Cartier test $H=H_1-H_2$
as a difference of two fixed globally nef classes gives
\[
 |-H^2\Lambda|
 \leq\sum_{a,b=1}^2\langle\gamma^{n-1}\rangle H_aH_b.
                                                        \tag{5}
\]
Only fixed nef tests enter this ordinary algebraic
positive-product comparison. Combined with finite ordered
$I\leq(n+1)(e(\xi)-e(\eta))$, it gives
\[
 |p_*\mu_\xi(h)-p_*\mu_\eta(h)|
 \leq C_{h,C}\sqrt{e(\xi)-e(\eta)}
 \quad(-C\leq\eta\leq\xi\leq0)                            \tag{6}
\]
for each fixed original rational model test $h$.
The proof and the actual fixed-twist generation are given in
whole201. The constant is uniform in both moving models.

\section{Two truncation lemmas}
For an integer $C>0$ put $\xi^C=\max(\xi,-C)$.
This is an actual generated model on the same $A$.

\begin{lemma}[Mass and energy of a finite truncation]\label{cap}
For $\xi\leq0$,
\[
 \|\mu_\xi-\mu_{\xi^C}\|_{\mathrm{TV}}
       \leq\frac{2(n+1)B_A(\xi)}{C}.                    \tag{7}
\]
If $\eta\leq\xi\leq0$ and both original costs are at most $D$, then
\[
 0\leq e_A(\xi^C)-e_A(\eta^C)
 \leq V_A|g_\xi-g_\eta|
        +K_nD^{1-\alpha}(e_A(\xi)-e_A(\eta))^\alpha .
                                                        \tag{8}
\]
\end{lemma}
\begin{proof}
The finite energy telescoping identity is
\[
 B_A(\xi)=\frac1{n+1}\sum_{a=0}^n
       \int(-\xi)\,\MA_A(\xi^a,0^{n-a}).
\]
Every summand is nonnegative, hence
$\int(-\xi)d\mu_\xi\leq(n+1)B_A(\xi)$.
Finite strict locality identifies $\mu_\xi$ and $\mu_{\xi^C}$
on $\{\xi>-C\}$. They have the same raw mass $V_A$, so the
masses of the complementary set agree. Their total variation
there is at most twice $\mu_\xi\{\xi\leq-C\}$.
Markov's inequality proves (7), including the equality level.
Pushforward can only decrease total variation.

The finite supporting plane at the lower potential gives
\[
 e_A(\xi^C)-e_A(\eta^C)
 \leq\int(\xi^C-\eta^C)\,d\mu_{\eta^C}.
\]
Pointwise $0\leq\xi^C-\eta^C\leq\xi-\eta$.
Monotonicity gives $B_A(\eta^C)\leq B_A(\eta)\leq D$.
Apply (4) with $\zeta=\eta^C$ to prove (8).
This requires no quadratic bound involving an unbounded model.
\end{proof}

\section{Exact meaning and existence of energy-complete joint limits}
Fix finitely many normalized $u_a\in\EE$ and, when needed, a
normalized relative test $v\in\EE$.
Choose their decreasing nonpositive model obstacle majorants
$f_{a,i}\downarrow u_a$ and $g_j\downarrow v$.
The obstacles and all their canonical envelopes are genuine
original functions; the physical full-system stages need not be
members of original $\EE$.

A joint net $\lambda$ has obstacle indices $i_a(\lambda)$,
a test index $j(\lambda)$ when used, and a nonempty cleared degree
$m(\lambda)$. All indices tend cofinally to their limits.
At that degree resolve all the complete lattices on one actual
$A_{m(\lambda)}$ and write
$\xi_{\lambda,a}=\psi_{m(\lambda),f_{a,i_a(\lambda)}}$ and
$\eta_\lambda=\psi_{m(\lambda),g_{j(\lambda)}}$.
The net is \emph{energy-complete} if, for each component and test,
\[
\begin{aligned}
 e_m(\xi_{\lambda,a})-E(P_\theta(f_{a,i_a(\lambda)}))&\longrightarrow0,
 &\xi_{\lambda,a}(v_0)&\longrightarrow0,\\
 e_m(\eta_\lambda)-E(P_\theta(g_{j(\lambda)}))&\longrightarrow0,
 &\eta_\lambda(v_0)&\longrightarrow0.
\end{aligned}                                           \tag{9}
\]
No weak or weighted output convergence is imposed in (9).
It follows from (S) that these auxiliary energies tend to the
actual original endpoint energies. On one cofinal tail all
nonpositive auxiliary costs are bounded by one finite $D$.

Such nets exist under (F).
For each finite tuple of obstacle indices and each tolerance
$\varepsilon>0$, its finite-list degree convergence gives
arbitrarily large common multiples with all errors in (9) below
$\varepsilon$. Index these valid tuples by increasing obstacle
indices, divisibility of degrees and decreasing tolerance.
Any finite set has a common successor: first choose dominating
obstacle indices and smaller tolerance, then a sufficiently large
common multiple satisfying the finite-list bounds.
All members beyond a tuple of tolerance $\varepsilon$ satisfy
their own errors at most $\varepsilon$, so (9) holds.
This construction also accommodates any fixed finite additional
model obstacles on common refinements.
There is no diagonal sequence or countability assumption.

\section{Every energy-complete diagonal measure}
\begin{theorem}[Unbounded diagonal identification]\label{diagonal}
Under the stated conditional inputs, every energy-complete net
for one $u\in\EE$ satisfies
\[
 p_{m*}\mu_{\xi_\lambda}\longrightarrow M(u)
                                                        \tag{10}
\]
weakly on the original $K$, independently of all degree selections
and common principalizations.
\end{theorem}
\begin{proof}
Fix one original model majorant index $k$ and an integer $C>0$.
Write $u_k=P_\theta(f_k)$ and $\xi_{m,k}=\psi_{m,f_k}$.
Once the moving index $i(\lambda)\geq k$, full filtration order gives
$\xi_\lambda\leq\xi_{m,k}$, on a common refinement.
For this fixed $k$, the upper auxiliary energy tends to $E(u_k)$;
the lower one tends to $E(u)$ by (9), and both generic values
tend to zero. Choose a single cost bound $D$ from the endpoint
energies; for each fixed $k$ it holds on a sufficiently late tail.
Equations (8) and (6) imply
\[
 \limsup_\lambda
 |p_{m*}\mu_{\xi_\lambda^C}(h)
       -p_{m*}\mu_{\xi_{m,k}^C}(h)|
 \leq C_{h,C}\sqrt{K_nD^{1-\alpha}
                        (E(u_k)-E(u))^\alpha}.          \tag{11}
\]
The generic term vanishes at this fixed $k$ before any $k$ limit.

The upper cap has physical stage at most
$w_k^C=\max(u_k,R-C)$, which is genuine and continuous on $U$.
Indeed $q_{T_m}\leq R$, so
$q_{T_m}+\max(\xi_{m,k},-C)\leq w_k^C$ on every lift.
Its quotient converges compact-uniformly to $w_k^C-R$,
and lies in $[-C,0]$ with generic error tending to zero.
It is an actual FS scalar after tensoring the horizontal factor.
Input (B) therefore identifies its pushed measure with $M(w_k^C)$.

Now let $k$ tend to its endpoint in (11).
Decreasing energy continuity makes its right side tend to zero,
and genuine decreasing weak continuity gives
$M(w_k^C)\to M(u^C)$ for $u^C=\max(u,R-C)$.
Thus $\mu_{\xi_\lambda^C}$ tends weakly to $M(u^C)$,
at this fixed $C$, on all model tests. Raw masses converge to $V$,
so uniform density passes all $h\in C(K)$.

By (7) the difference from the untruncated output is at most
$2(n+1)D\|h\|_\infty/C$, uniformly on the energy-complete tail.
Finally $u^C\downarrow u$ genuinely as $C\to\infty$,
so $M(u^C)\to M(u)$.
Take the moving limit at fixed $C$ before sending $C$ to infinity.
This proves (10); it uses no uniform quadratic estimate for
the original unbounded $\xi_\lambda$.
\end{proof}

\section{All mixed slots, without extra barycenter hypotheses}
\begin{theorem}[Mixed and energy coefficient identification]\label{mixed}
For an energy-complete net of finitely many inputs, every finite
auxiliary mixed measure converges weakly to the corresponding
genuine mixed measure. The energy polynomials converge
coefficientwise to their genuine counterparts.
Individual endpoint completeness (9) suffices.
\end{theorem}
\begin{proof}
Fix a rational simplex vector $t=(t_a)$ and put
$\xi_{\lambda,t}=\sum_a t_a\xi_{\lambda,a}$.
At fixed upper indices $k=(k_a)$ put
$\xi_{m,k,t}=\sum_a t_a\psi_{m,f_{a,k_a}}$ and
$w_{k,t}=\sum_a t_aP_\theta(f_{a,k_a})$.
These are actual rational tensor systems on the \emph{same} $A_m$.
After clearing powers their physical scalar is
$q_{T_m}+\xi_{\lambda,t}$: the horizontal factor is $T_m$,
since $\sum_a t_a=1$, not a new full-system factor in a higher degree.
The moving lower scalar is at most the fixed upper scalar
once every input index dominates $k$.
Concavity bounds its cost by $D$ and its generic value tends to zero.

The finite supporting plane at $\xi_{\lambda,t}$ and (4) give
\[
\begin{split}
0&\leq e_m(\xi_{m,k,t})-e_m(\xi_{\lambda,t})\\
&\leq\sum_a t_a\left[
 V_m|g_{m,k_a}-g_{\lambda,a}|
 +K_nD^{1-\alpha}
  (e_m(\psi_{m,f_{a,k_a}})-e_m(\xi_{\lambda,a}))^\alpha
 \right].
\end{split}                                             \tag{12}
\]
For fixed $k$, (B) identifies the upper energy with $E(w_{k,t})$.
First take the moving limit in (12), and then $k$ to its endpoint.
The right side tends to zero, while
$w_{k,t}\downarrow w_t=\sum_a t_a u_a$ with decreasing energies.
Consequently $e_m(\xi_{\lambda,t})\to E(w_t)$.
Thus energy completeness for these barycenters is a conclusion.

The proof of Theorem~\ref{diagonal} applies verbatim to this
ordered lower/upper pair: fixed upper barycenters and their caps
satisfy (B), all costs are bounded, and (12) supplies the required
energy limit. Hence
$p_{m*}\MA(\xi_{\lambda,t})\to M(w_t)$.

Choose one finite rational unisolvent set on the simplex,
before taking any limit: degree $n$ for the raw measure polynomial
and degree $n+1$ for the energy polynomial.
Finite intersection multilinearity holds on the actual nef $A_m$,
even when it is not ample; it does not concern multilinearity of
big positive products.
Invert the fixed finite evaluation matrices.
The preceding diagonal and energy limits identify every coefficient
with the genuine coefficients constructed in whole202.
Each mixed coefficient remains a positive measure of raw mass
$V_m\to V$; finite weak convergence survives the fixed linear inverse.
The same coefficient argument gives every real-weight polynomial
identity, without an exchange of a varying weight and a moving limit.
\end{proof}

\section{Simultaneously changing paired relative tests}
\begin{theorem}[Paired moving weighted convergence]\label{weighted}
Add an independently changing canonical test index for
$v\in\EE$ satisfying (9) on the same actual $A_m$.
If $\sigma_\lambda$ is any $n$-slot raw mixed auxiliary output
of the input models, then for every $h\in C(K)$,
\[
 \int (h\circ p_m)(-\eta_\lambda)\,d\sigma_\lambda
 \longrightarrow
 \int_U h(R-v)\,dM(u_1,\ldots,u_n).                     \tag{13}
\]
All limits are full directed joint limits; inputs and test indices
may change independently. In particular the unweighted moment
and the signed $h$-weighted moment have the same choice-independent
limits.
\end{theorem}
\begin{proof}
Finite positive polarization gives
$0\leq\sigma_\lambda\leq(n^n/n!)\MA(\bar\xi_\lambda)$,
where $\bar\xi_\lambda=n^{-1}\sum_a\xi_{\lambda,a}$.
Its cost is at most $D$ by concavity.
Fix one test upper index $k$ and write
$v_k=P_\theta(g_k)$, $\eta_{m,k}=\psi_{m,g_k}$.
Full filtration order gives $\eta_\lambda\leq\eta_{m,k}$
eventually. Apply (4) at the average output:
\[
\begin{split}
\limsup_\lambda
 \int(\eta_{m,k}-\eta_\lambda)\,d\sigma_\lambda
 &\leq\frac{n^n}{n!}K_nD^{1-\alpha}
                      (E(v_k)-E(v))^\alpha .
\end{split}                                             \tag{14}
\]
Generic errors vanish at fixed $k$ and are not discarded earlier.
The right side tends to zero as $k$ tends to its endpoint.

At this fixed $k$, $\eta_{m,k}$ is bounded and converges
uniformly on compact subsets of $U$ to $v_k-R$.
Theorem~\ref{mixed} gives the weak mixed output limit
$\sigma=M(u_1,\ldots,u_n)$; by (1) it gives $Z$ zero mass.
A fixed original continuous cutoff supported in $U$ can have
omitted $\sigma$-mass as small as desired.
Weak convergence controls the same omitted moving mass.
On its compact support the changing bounded test converges uniformly.
First choose this collar, then take the moving limit, then remove
the collar. It follows that
\[
 \int(h\circ p_m)(-\eta_{m,k})d\sigma_\lambda
       \longrightarrow\int_Uh(R-v_k)d\sigma .
\]
No unbounded test is passed by weak convergence.
Finally genuine ordered absolute errors, or the canonical weighted
continuity in (M), give
$\int(v_k-v)d\sigma\to0$.
Multiply (14) by $\|h\|_\infty$ and send $k$ to its endpoint.
This proves (13).
\end{proof}

\section{The singular-reference distinction and exact remaining scope}
Equation (13) concerns the paired auxiliary quotient
$-\eta_\lambda$. At every divisorial atom it is exactly
the difference between the physical reference $q_{T_m}$
and the paired physical test $q_{T_m}+\eta_\lambda$.
This statement uses no original finite-energy membership for
either physical section function.

It must not be restated as
\[
 \int_K h(R-v)\,d(p_{m*}\sigma_\lambda)
       \longrightarrow\int_Kh(R-v)d\sigma
                                                        \tag{15}
\]
for an arbitrary fixed genuine unbounded $v$.
Weak convergence and (13) alone do not control its difference
from a physical test. At finite divisorial atoms that difference
contains the nonnegative singular-reference defect
$R-q_{T_m}$ and the error between a genuine $v$ and its
physical minorant. Neither is bounded by (9) in this paper.
At supported poles one must not subtract two infinite values.
Thus (15), and a uniform absolute original-test approximation
for every physical auxiliary output, remain separate assertions.

The original genuine operator does have its canonical weighted
continuity under (M). This paper additionally identifies its
energy-complete full-system diagonal and mixed weak limits
and its \emph{paired moving} relative weighted moments.
It proves the needed auxiliary ordered-tail estimates by (4);
there is no circular output convergence in the admissibility
condition (9), no bound on growing model sup norms, and no
uniform quadratic estimate involving unbounded slots.
It does not claim arbitrary degree paths with uncontrolled
energy error, nor a proof of the actual nonample antecedents.
Further nonzero zero-volume theory is outside the positive-volume
argument. Original P16.5 remains open and no acceptance count changes.

\section{Source application and reproducibility scope}
The whole finite and conditional proofs in159,201 and202 were
read. Formula (3) is precisely159's finite absolute moment;
the actual fixed-twist quadratic estimate (5)--(6) is201's
bounded finite-model argument. Input (B) is202's bounded physical
identification, including its separate original and auxiliary
collars, full filtered-quotient missing-volume derivative and
generic normalization terms. These are unreviewed proof bodies,
not independently accepted theorems.
The new arguments supplied here are (7)--(8), the energy-complete
directed filter, (11)--(12), and the paired weighted passage
(14). The finite identities and conditional proofs are written
above; compilation does not certify them.
No new primary-source retrieval, whole-primary reading,
recursive-background proof or original PDF-byte hash is claimed.
The ordinary algebraic comparisons in (5) retain the exact
smooth characteristic-zero hypotheses of BFJ and are used
only in the bounded finite range established in201.
The existing user editor remains unchanged.
\begin{thebibliography}{9}
\bibitem{BJ} S. Boucksom and M. Jonsson,
\emph{Global pluripotential theory over a trivially valued field},
\href{https://arxiv.org/pdf/1801.08229v4}{arXiv:1801.08229v4}.
\bibitem{BFJ} S. Boucksom, C. Favre and M. Jonsson,
\emph{Differentiability of volumes of divisors and a problem of Teissier},
\href{https://arxiv.org/html/math/0608260v2}{arXiv:math/0608260v2}.
\end{thebibliography}
\end{document}
